%@ pages 144-146
%@ status translated
% Original section 14, "Limites", pp. 144--146.

\clearpage
\section{Limits}
\label{limitespage}

In this section we recall a few facts concerning limits. In \cite{limits} one
defines the notions of limit and colimit in an $n$-category. One can give the
same definitions, {\em mutatis mutandis}, in a Segal $n$-category. We refer the
reader to \cite{limits}, 3.1 and 3.2 for these definitions. Transcribing the
proofs of \cite{limits} into the framework of Segal $n$-categories, one
obtains:

\begin{theorem}
The Segal $(n+1)$-category $nSeCAT'$ admits all limits and colimits indexed by
(small) Segal $(n+1)$-categories.
\end{theorem}

Let us recall how to compute a limit concretely (this is the construction used
in \cite{limits} 4.1). If $\Yy$ is an $(n+1)$-category equipped with a
morphism $F: \Yy \rightarrow nSeCAT'$, then one has (\cite{limits} 4.1.2)
$$
\lim _{\leftarrow , \Yy} F
= \underline{Hom}(\Yy , nSeCAT')_{1/}(\ast _{\Yy}, F).
$$

On the other hand, one has at one's disposal the notions of homotopy limit and
colimit ($holim$ and $hocolim$) for diagrams (indexed by a small
$1$-category) in an essentially arbitrary cmc. These notions are due to
Bousfield and Kan \cite{BousfieldKan}; see also
%%%%%% added in v3 : %%%%%%
Vogt \cite{Vogt}, Edwards and Hastings \cite{EH},
Cordier and Porter \cite{CordierPorter2},
%%%%%%%%%%
Hirschhorn \cite{Hirschhorn},
Dwyer--Hirschhorn--Kan \cite{DHK}.

Normally we shall use the notation $holim$ for constructions of
Bousfield--Kan type, and the notation $\lim$ for limits in the sense of
\cite{limits}. Fortunately, the two notions coincide for limits of Segal
$n$-categories indexed by $1$-categories. This is essentially a consequence of
the following proposition. We leave it to the reader to make the analogous
comparison with the $holim$ of \cite{Hirschhorn} for the cmc $nSePC$.

\begin{proposition}
\label{calclim}
Suppose that $\Yy$ is a category and $F$ a Segal $n$-prestack over $\Yy$,
which one can also regard as a morphism $\Yy ^o\rightarrow nSeCAT$. Let
$F\rightarrow F'$ be an object-by-object weak equivalence to a Segal
$n$-prestack that is fibrant for the coarse topology. Then one has
equivalences of Segal $n$-categories
$$
\Gamma (\Yy , F')
\stackrel{\cong}{\rightarrow}
$$
$$
\lim _{\leftarrow , \Yy ^o} F' \cong
\lim _{\leftarrow , \Yy ^o} F .
$$
Here the limits are those defined following \cite{limits}.
\end{proposition}
{\em Proof:}
One has the equality
$\Gamma (\Yy , F')=\Gamma \underline{Hom}(\ast _{\Yy}, F')$ because the two
sides represent the same functor on the category of Segal $n$-precategories.
In particular, one has
$$
\Gamma (\Yy , F')= nSeCHAMP (\Yy )_{1/}(\ast, F').
$$
On the other hand, by (\cite{limits} 4.0.1 and 4.1.2), one has
$$
\lim _{\leftarrow , \Yy ^o} F' = \underline{Hom}(\Yy ^o, nSeCAT')_{1/} (\ast
_{\Yy}, F').
$$
Now the fact that the functor $\Phi$ of Theorem \ref{correlation} is fully
faithful (cf.\ Proposition \ref{pleinfidele}) means exactly that
$$
nSeCHAMP (\Yy )_{1/}(\ast, F')\rightarrow
\underline{Hom}(\Yy ^o, nSeCAT')_{1/} (\ast
_{\Yy}, F')
$$
is an equivalence. Finally the equivalence between the limits of $F$ and $F'$
follows from the invariance properties of the notion of limit
\cite{limits}.
\eop

There is a similar characterisation for colimits, which uses the HBKQ-type
structure (\S5). We state it below although we shall not need it. This
proposition can also be regarded as a compatibility statement between the
definition of $hocolim$ of Bousfield--Kan and the construction of
\cite{limits}.

\begin{proposition}
\label{calccolim}
Let $\Yy$ be a category and $F$ a Segal $n$-prestack over $\Yy$, also regarded
as a morphism $\Yy ^o\rightarrow nSeCAT$. Let $F'\rightarrow F$ be an
object-by-object weak equivalence with $F'$ cofibrant for the cmc structure of
HBKQ type, for the coarse topology. Then one has equivalences of Segal
$n$-categories
$$
colim_{nSePC}(F')
\stackrel{\cong}{\rightarrow}
$$
$$
\lim _{\rightarrow , \Yy ^o} F' \stackrel{\cong}{\rightarrow}
\lim _{\rightarrow , \Yy ^o} F
$$
where $colim_{nSePC}(F')$ is the standard colimit of the functor
$\Yy \rightarrow nSePC$ from $\Yy$ to the category of Segal $n$-precategories,
while the other colimits are the colimits defined following \cite{limits}.
\end{proposition}
\eop

Thanks to Proposition \ref{calclim} one obtains a version of our criteria
which uses the limits of \cite{limits}. This version of the definition of
stacks was suggested in \cite{limits}.

\begin{corollary}
\label{critereaveclim}
Let $A$ be a Segal $n$-prestack on $\Xx$ whose values are Segal
$(n+1)$-categories, corresponding to a morphism $\Xx \rightarrow nSeCAT$.
Then for $A$ to be a Segal $n$-stack it is necessary and sufficient that:
\newline
(a)\,\, for every $X\in \Xx$ and every $x,y\in A_0(X)$ the Segal
$(n-1)$-prestack $A_{1/}(x,y)$ on $\Xx /X$ be a Segal $(n-1)$-stack
(resp.\ $Path^{x,y}A$ be a Segal $0$-stack, in the case $n=0$); and
\newline
(b)\,\, for every $X\in \Xx$ and every sieve $\Bb \subset \Xx /X$ of the
topology, the morphism
$$
A(X)\rightarrow \lim _{\leftarrow , \Bb ^o} A|_{\Bb}
$$
be essentially surjective.

%%%% the following is changed in v3 %%%%%%
Or again, for a (non-Segal) $n$-prestack $A$ to be an $n$-stack it is
necessary and sufficient that:
\newline
(c)\,\, for every $X\in \Xx$ and every sieve $\Bb \subset \Xx /X$ of the
topology, the morphism
$$
A(X)\rightarrow \lim _{\leftarrow , \Bb ^o} A|_{\Bb}
$$
be an equivalence of Segal $n$-categories.
\end{corollary}
{\em Proof:}
In view of Proposition \ref{calclim}, condition (b) is equivalent to that of
\ref{critere} or \ref{0critere} (the conditions (a) are identical). Likewise
condition (c) is equivalent to that of \ref{newcritere} (which, according to
the {\tt v3} correction, works only for non-Segal $n$-prestacks).
\eop

{\em Proof of the second part of Theorem \ref{correlation}:} if the essential
surjectivity of $\Phi$ is known for the first part of \ref{correlation} (which
will be proved in \ref{strictif3} below), the second part of \ref{correlation}
now follows from condition (c) of the corollary above.

\bigskip
